\subsection{Выпуклый симплексный метод Зангвилла}
\label{subsec:zangwill}

\[
  \varphi(x) \to \min
  \qquad\text{при}\qquad
  Ax = b , \quad x \ge 0 ,
\]
где $\varphi$ — дифференцируемая (желательно выпуклая) функция, $A$ — матрица
$m \times n$. Допустимое множество $S = \{\, x \mid Ax = b,\ x \ge 0 \,\}$.
Метод Зангвилла — обобщение симплекс-метода ЛП на случай нелинейной выпуклой
цели.

\subsubsection*{Линейная подзадача}

На $k$-й итерации $\varphi$ заменяется линейной аппроксимацией в точке $x^k$:
\[
  \varphi(x) \approx \varphi(x^k) + \nabla\T\varphi(x^k)\,(x - x^k) .
\]
Слагаемые $\varphi(x^k)$ и $-\nabla\T\varphi(x^k)x^k$ от $x$ не зависят,
поэтому минимизация аппроксимации сводится к задаче
\[
  g_k\T x \to \min_{x \in S} ,
  \qquad g_k = \nabla\varphi(x^k) ,
\]
то есть к задаче \emph{линейного программирования} — отсюда и название
«симплексный». Её решение обозначим $y^k$, направление движения
\[
  d^k = y^k - x^k .
\]

\begin{theorem}\label{thm:zangwill-descent}
Если $g_k\T\big(y^k - x^k\big) < 0$, то $d^k = y^k - x^k$ — направление спуска.
\end{theorem}

\begin{proof}
$\nabla\T\varphi(x^k)d^k = g_k\T\big(y^k - x^k\big) < 0$ — это и есть
определение направления спуска.
\end{proof}

\begin{theorem}\label{thm:zangwill-opt}
Пусть $\varphi$ выпукла и дифференцируема, $S$ выпукло, $y^k$ решает
$g_k\T x \to \min_{x\in S}$. Если $g_k\T\big(y^k - x^k\big) = 0$, то $x^k$ —
решение задачи $\varphi(x) \to \min_{x \in S}$.
\end{theorem}

\begin{proof}
Так как $y^k$ минимизирует линейную функцию на $S$, для любого $x \in S$
выполнено $g_k\T y^k \le g_k\T x$. Вычитая $g_k\T x^k$,
\[
  \underbrace{g_k\T\big(y^k - x^k\big)}_{=\ 0} \ \le\ g_k\T\big(x - x^k\big)
  \qquad\Longrightarrow\qquad
  \nabla\T\varphi(x^k)\big(x - x^k\big) \ \ge\ 0 \quad \forall x \in S .
\]
Это условие экстремума выпуклой функции на выпуклом множестве; в силу
выпуклости $\varphi$ оно достаточно, значит $x^k$ — глобальный минимум.
\end{proof}

\begin{theorem}\label{thm:zangwill-feasible}
Если $x^k \in S$, $y^k \in S$ и $S$ выпукло, то
$x^{k+1} = x^k + \alpha_k\big(y^k - x^k\big) \in S$ при любом
$\alpha_k \in [0,1]$.
\end{theorem}

\begin{proof}
$x^{k+1} = (1-\alpha_k)x^k + \alpha_k y^k$ — выпуклая комбинация двух точек
$S$; по выпуклости она лежит в $S$.
\end{proof}

Шаг выбирается одномерной минимизацией
\[
  \alpha_k = \argmin_{0 \le \alpha \le 1} \varphi\big(x^k + \alpha d^k\big) .
\]

\begin{algo}[title={Выпуклый симплексный метод Зангвилла}]
\begin{steps}
  \item Задать допустимую $x^0 \in S$, точность $\eps > 0$, $k = 0$.
  \item Вычислить $g_k = \nabla\varphi(x^k)$.
  \item Решить линейную подзадачу $y^k = \argmin_{x\in S} g_k\T x$ (симплекс-метод).
  \item Построить $d^k = y^k - x^k$.
  \item Если $g_k\T d^k \ge -\eps$ — остановиться; в выпуклом случае $x^k$
        принимается за приближённое решение.
  \item Выбрать $\alpha_k = \argmin_{0\le\alpha\le1}\varphi(x^k + \alpha d^k)$.
  \item $x^{k+1} = x^k + \alpha_k d^k$, $k := k+1$, перейти к шагу~2.
\end{steps}
\end{algo}

\subsubsection*{Классический вариант: смена одной небазисной переменной}

\begin{idea}[title={Откуда этот материал}]
Выше изложен вариант со вспомогательной задачей ЛП на всём $S$. В
экзаменационной формулировке под методом Зангвилла понимается \emph{классический}
выпуклый симплексный метод, в котором задача ЛП не решается вовсе: по правилу
$u$/$v$ выбирается ровно одна небазисная переменная, и шаг делается вдоль
одного ребра — ровно как в симплекс-методе ЛП. Он и разобран ниже.
\end{idea}

Пусть $x = (x_B, x_N)$, $A = (B,N)$, $B$ невырождена, $x_B > 0$. Как и в методе
Вульфа (\S\ref{subsec:reduced-gradient}), вводится приведённый градиент
\[
  r_j = \frac{\partial\varphi}{\partial x_j}
      - \nabla_B\T\varphi\,B^{-1}A_j , \qquad j \in N ,
\]
и \textbf{критерий оптимальности} (условия Куна--Таккера):
\[
  r_j \ge 0 \ \text{при } x_j = 0 ,
  \qquad
  r_j = 0 \ \text{при } x_j > 0 .
\]

\paragraph{Выбор переменной.}
\[
  u = \min_{j \in N} r_j = r_s ,
  \qquad
  v = \max_{j \in N} x_j r_j = x_t r_t .
\]
Останов при $u \ge 0$ и $v = 0$ — это в точности критерий выше. Иначе:
\begin{gather*}
  \abs{u} \ge v \ \Longrightarrow\ \text{увеличиваем } x_s \
    (\sigma = +1,\ q = s) ; \\
  \abs{u} < v \ \Longrightarrow\ \text{уменьшаем } x_t \
    (\sigma = -1,\ q = t) .
\end{gather*}

\paragraph{Направление.}
$d_q = \sigma$, остальные небазисные нули, $d_B = -\sigma B^{-1}A_q$. Тогда
$Ad = \sigma A_q - \sigma A_q = 0$, то есть равенства сохраняются, и
$\nabla\T\varphi\,d = r_q d_q$:
\begin{itemize}[nosep,left=0pt]
  \item $\sigma = +1$: $r_q = u < 0$, значит $\nabla\T\varphi\,d < 0$;
        увеличение $x_q \ge 0$ допустимо всегда;
  \item $\sigma = -1$: $r_q > 0$ и $x_q > 0$, значит
        $\nabla\T\varphi\,d = -r_q < 0$; уменьшать $x_q$ можно до нуля.
\end{itemize}
Базисные компоненты строго положительны, поэтому при малом шаге точка остаётся
допустимой, а максимальный шаг —
$\alpha_{\max} = \min\{-x_i/d_i : d_i < 0\}$ (при $\sigma = -1$ учитывается и
$x_q$).

\begin{theorem}\label{thm:zangwill-kkt}
Предельные точки последовательности $\{x^k\}$ — точки Куна--Таккера; при
выпуклой $\varphi$ они являются решениями задачи.
\end{theorem}

\begin{idea}[title={Доказательство (схема)}]
$\varphi(x^k)$ убывает и ограничена снизу, значит сходится. Базисов конечное
число, поэтому можно выделить подпоследовательность с фиксированным базисом;
при фиксированном базисе направление $d(x)$ непрерывно по $x$. Если бы для
предельной точки $\bar{x}$ было $d(\bar{x}) \ne 0$, нашёлся бы шаг
$\bar\alpha > 0$ с $\varphi(\bar{x} + \bar\alpha d(\bar{x})) < \varphi(\bar{x})$,
что противоречит сходимости $\varphi(x^k)$. Значит $d(\bar{x}) = 0$, то есть
$u \ge 0$ и $v = 0$.
\end{idea}

\begin{idea}[title={Чем это отличается от метода условного градиента}]
Вариант со вспомогательной задачей ЛП — это тот же метод Франка--Вульфа
(\S\ref{subsec:cond-gradient}), и теоремы~\ref{thm:zangwill-descent}--\ref{thm:zangwill-feasible}
дословно повторяют результаты того параграфа. Разница в акценте:
\begin{itemize}[nosep,left=0pt]
  \item в \S\ref{subsec:cond-gradient} множество $S$ — произвольный выпуклый
        компакт, и подзадача решается «как получится»;
  \item у Зангвилла $S$ задано как $Ax = b$, $x \ge 0$, подзадача — честная
        задача ЛП в канонической форме, и решается она симплекс-методом с
        тёплым стартом от предыдущего базиса.
\end{itemize}
Поэтому метод Зангвилла — это практическая реализация условного градиента для
линейных ограничений, где вся техника §1--§4 работает внутри каждой итерации.
Скорость та же сублинейная, $O(1/k)$.
\end{idea}

\begin{example}
$\varphi(x) = (x_1-0{,}2)^2 + (x_2-0{,}3)^2 + (x_3-0{,}5)^2 \to \min$ при
$x_1+x_2+x_3 = 1$, $x \ge 0$; старт $x^0 = (1;\,0;\,0)\T$,
$\varphi(x^0) = 0{,}98$.

Итерация 1: $g_0 = (1{,}6;\,-0{,}6;\,-1)\T$. Линейная подзадача
$\min(1{,}6x_1 - 0{,}6x_2 - x_3)$ на симплексе достигает минимума в вершине с
наименьшим коэффициентом, то есть $y^0 = (0;\,0;\,1)\T$. Направление
$d^0 = (-1;\,0;\,1)\T$, $g_0\T d^0 = -2{,}6 < 0$. Одномерная минимизация
$\varphi = 0{,}98 - 2{,}6\alpha + 2\alpha^2$ даёт $\alpha_0 = 0{,}65$, откуда
$x^1 = (0{,}35;\,0;\,0{,}65)\T$, $\varphi = 0{,}135$.

Итерация 2: $g_1 = (0{,}3;\,-0{,}6;\,0{,}3)\T$, минимальный коэффициент при
$x_2$, значит $y^1 = (0;\,1;\,0)\T$, $d^1 = (-0{,}35;\,1;\,-0{,}65)\T$,
$g_1\T d^1 = -0{,}9 < 0$. Минимизация
$\varphi = 0{,}135 - 0{,}9\alpha + 1{,}545\alpha^2$ даёт
$\alpha_1 = 0{,}9/3{,}09 \approx 0{,}2913$, откуда
$x^2 \approx (0{,}248;\,0{,}291;\,0{,}461)\T$, $\varphi \approx 0{,}0039$.
Итого $0{,}98 \to 0{,}135 \to 0{,}0039$.
\end{example}

\begin{idea}
Видно характерное для метода поведение: решение линейной подзадачи всегда
попадает в \emph{вершину} симплекса, и без одномерного поиска точка прыгала бы
между вершинами. Именно шаг $\alpha_k \in [0,1]$ позволяет приближаться к
внутреннему оптимуму $x^* = (0{,}2;\,0{,}3;\,0{,}5)\T$ как к выпуклой
комбинации вершин.
\end{idea}
